% Vietnamese translation for OI Public Library
% Source-PDF: IOI中国国家候选队论文/2001/毛子青.pdf
\documentclass[11pt]{article}
\usepackage{oipl-vn}

\title{Kỹ thuật tối ưu hóa thuật toán quy hoạch động}
\author{Mao Ziqing}
\date{}

\begin{document}
\maketitle

\origtitle{Optimization techniques for dynamic programming algorithms}
\sourcepath{IOI China candidate team papers / 2001 / Mao Ziqing.pdf}

\paragraph{Từ khóa.}
Quy hoạch động; độ phức tạp thời gian; tối ưu hóa; trạng thái.

\section*{Tóm Tắt}

Quy hoạch động là một phương pháp thiết kế chương trình thường dùng trong các
kỳ thi tin học. Bài viết này tập trung thảo luận việc tối ưu hóa hiệu suất
thời gian khi dùng tư tưởng quy hoạch động để giải bài. Toàn văn gồm bốn
phần: trước hết bàn về tính khả thi và sự cần thiết của tối ưu hóa thời gian
trong quy hoạch động; tiếp đó nêu các yếu tố quyết định độ phức tạp thời gian
của quy hoạch động; sau đó lần lượt trình bày phương pháp tối ưu hóa từng yếu
tố; cuối cùng là phần kết luận.

\tableofcontents

\section{Dẫn nhập}

Quy hoạch động là một phương pháp thiết kế chương trình quan trọng, có phạm vi
ứng dụng rộng trong các kỳ thi tin học.

Khi dùng quy hoạch động để giải bài, khá nhiều bài có đặc điểm là tốn nhiều
bộ nhớ nhưng hiệu suất thời gian cao. Vì vậy, khi nghiên cứu cách giải bằng
quy hoạch động, người ta thường chú ý nhiều hơn đến việc tối ưu hóa độ phức
tạp không gian, dùng nhiều kỹ thuật khác nhau để khống chế nhu cầu bộ nhớ
trong phạm vi mà phần cứng và phần mềm có thể chịu được. Tuy nhiên, cũng có
một số bài mà khi dùng tư tưởng quy hoạch động, hiệu suất thời gian vẫn chưa
đáp ứng yêu cầu, trong khi thuật toán còn khả năng cải tiến. Khi đó cần xét
đến tối ưu hóa hiệu suất thời gian.

Bài viết này thảo luận việc tối ưu hóa lời giải quy hoạch động đã có, trong
điều kiện đã xác định sẽ dùng tư tưởng quy hoạch động để giải bài, nhằm hạ độ
phức tạp thời gian của thuật toán và giúp thuật toán áp dụng được cho quy mô
lớn hơn.

\section{Phân tích độ phức tạp thời gian của quy hoạch động}

Khi dùng quy hoạch động để giải bài, lý do quan trọng khiến nhiều bài có hiệu
suất thời gian cao là nó giảm bớt phần ``dư thừa''. Ở đây, ``dư thừa'' chỉ
những phép tính không cần thiết hoặc những phần bị tính lặp lại; mức độ dư
thừa của thuật toán là yếu tố then chốt quyết định hiệu suất. Trong khi liên
tục thu nhỏ quy mô bài toán, quy hoạch động ghi lại nghiệm của các bài toán
con đã giải, tận dụng đầy đủ các kết quả ấy, tránh hiện tượng phải giải đi
giải lại cùng một bài toán con, từ đó giảm dư thừa.

Tuy vậy, khi quy hoạch động giải bài vẫn còn dư thừa. Nó chủ yếu gồm: giải
những bài toán con vô dụng, tham chiếu những kết quả không có ý nghĩa đối với
kết quả cuối cùng, v.v.

Ta có thể mô tả trực quan các yếu tố quyết định độ phức tạp thời gian của quy
hoạch động như sau:
\[
\begin{aligned}
T &= S\cdot C\cdot U,\\
S &= \text{tổng số trạng thái},\\
C &= \text{số trạng thái chuyển của mỗi trạng thái},\\
U &= \text{thời gian cho mỗi lần chuyển trạng thái}.
\end{aligned}
\]
Các phần sau lần lượt thảo luận cách tối ưu hóa ba yếu tố này. Cần chỉ ra
rằng ba yếu tố không độc lập với nhau, mà liên hệ, mâu thuẫn và thống nhất
với nhau. Có lúc tối ưu được một yếu tố thì hai yếu tố còn lại cũng được cải
thiện theo; có lúc tối ưu một yếu tố lại phải trả giá bằng việc tăng một yếu
tố khác. Vì vậy, khi tối ưu hóa cần giữ cái nhìn toàn cục và đạt cân bằng
giữa ba yếu tố.

\section{Tối ưu hóa hiệu suất thời gian của quy hoạch động}

\subsection{Giảm tổng số trạng thái}

Ta biết rằng quá trình giải bằng quy hoạch động thực chất là quá trình tính
giá trị của tất cả trạng thái; do đó quy mô trạng thái ảnh hưởng trực tiếp
đến hiệu suất thời gian của thuật toán. Vì vậy, giảm tổng số trạng thái là
một phần quan trọng trong tối ưu hóa quy hoạch động. Mục này thảo luận một số
phương pháp giảm tổng số trạng thái.

\subsubsection{Cải tiến cách biểu diễn trạng thái}

Quy mô trạng thái có quan hệ mật thiết với cách biểu diễn trạng thái. Giảm
tổng số trạng thái bằng cách cải tiến biểu diễn trạng thái là một phương pháp
được dùng khá phổ biến.

\paragraph{Ví dụ 1: nhóm nhạc Raucous Rockers (USACO 1996).}

\paragraph{Mô tả bài toán.}
Có \(n\) bài hát quý do nhóm Raucous Rockers thu âm. Ta dự định chọn một số
bài trong đó để phát hành trên \(m\) đĩa, mỗi đĩa chứa nhiều nhất \(t\) phút
nhạc, và các bài hát trên đĩa không được chồng lấn nhau. Việc chọn bài phải
theo các tiêu chuẩn sau:
\begin{enumerate}
  \item Các bài hát trong bộ đĩa phải được sắp theo thứ tự sáng tác.
  \item Tổng số bài hát được chứa là nhiều nhất có thể.
\end{enumerate}
Cho \(n,m,t\) và độ dài của \(n\) bài hát theo thứ tự sáng tác. Không bài nào
dài hơn một đĩa, và không thể đặt tất cả bài hát lên các đĩa. Hãy xuất số bài
hát nhiều nhất có thể chứa được, với
\[
  1\le n,m,t\le 20.
\]

\paragraph{Phân tích thuật toán.}
Điều kiện các bài hát trong đĩa phải theo thứ tự sáng tác thỏa yêu cầu không
hậu hiệu của quy hoạch động, gợi ý ta dùng quy hoạch động.

Phân tích cho thấy bài toán có tính chất cấu trúc con tối ưu: giả sử trong
phương án thu âm tối ưu, bài hát thứ \(i\) được ghi từ phút thứ \(k\) của đĩa
thứ \(j\). Khi đó \(j-1\) đĩa trước và \(k-1\) phút đầu của đĩa thứ \(j\) là
phương án thu âm tối ưu cho các bài \(1..i-1\). Nói cách khác, nghiệm tối ưu
của bài toán chứa nghiệm tối ưu của bài toán con.

Gọi độ dài của \(n\) bài hát sau khi sắp theo thứ tự sáng tác là
\(\operatorname{long}[1..n]\). Một cách biểu diễn trạng thái là
\[
  g[i,j,k]\quad (0\le i\le n,\ 0\le j\le m,\ 0\le k<t),
\]
biểu thị số bài hát nhiều nhất có thể ghi khi xét \(i\) bài đầu, dùng \(j\)
đĩa cộng thêm \(k\) phút. Nghiệm tối ưu của bài toán là \(g[n,m,0]\). Vì bài
hát \(i\) có hai khả năng, phát hành hoặc không phát hành, và còn phải xét
\(k\) phút cộng thêm có đủ ghi bài \(i\) hay không, ta có các công thức
chuyển trạng thái và điều kiện biên sau.

Khi \(k\ge \operatorname{long}[i]\) và \(i\ge1\):
\[
  g[i,j,k]=\max\{g[i-1,j,k-\operatorname{long}[i]]+1,\ g[i-1,j,k]\}.
\]
Khi \(k<\operatorname{long}[i]\) và \(i\ge1\):
\[
  g[i,j,k]=
  \max\{g[i-1,j-1,t-\operatorname{long}[i]]+1,\ g[i-1,j,k]\}.
\]
Điều kiện biên là
\[
  g[0,j,k]=0 \qquad (0\le j\le m,\ 0\le k<t).
\]

Thuật toán trên có số trạng thái \(O(nmt)\), số trạng thái chuyển cho mỗi
trạng thái là \(O(1)\), thời gian cho mỗi lần chuyển là \(O(1)\), nên tổng
thời gian là \(O(nmt)\). Với \(n,m,t\le20\), thuật toán đáp ứng yêu cầu.

\paragraph{Tối ưu hóa thuật toán.}
Khi quy mô dữ liệu lớn hơn, thuật toán trên không còn đáp ứng yêu cầu. Ta xét
cách tăng hiệu suất thời gian bằng cách cải tiến biểu diễn trạng thái.

Mục tiêu tối ưu của bài là dùng một số đĩa có độ dài cho trước để ghi càng
nhiều bài càng tốt. Điều này thực chất tương đương với: khi cần ghi một số
lượng bài hát cho trước, hãy dùng ít đĩa nhất có thể. ``Ít đĩa nhất'' nghĩa
là số đĩa đầy đủ đã dùng là ít nhất, hoặc khi số đĩa đầy đủ như nhau thì số
phút cộng thêm là ít nhất. Dưới mục tiêu tối ưu ấy, bài toán cũng có tính
cấu trúc con tối ưu. Giả sử \(D\) là phương án dùng ít đĩa nhất để chọn \(j\)
bài trong \(i\) bài đầu. Nếu phương án chọn bài thứ \(i\), thì \(D-\{i\}\) là
phương án dùng ít đĩa nhất để chọn \(j-1\) bài trong \(i-1\) bài đầu; nếu
không chọn bài thứ \(i\), thì \(D\) là phương án dùng ít đĩa nhất để chọn
\(j\) bài trong \(i-1\) bài đầu. Nghiệm tối ưu của bài toán vẫn chứa nghiệm
tối ưu của bài toán con.

Biểu diễn trạng thái cải tiến là
\[
  g[i,j]=(a,b),\quad
  0\le i\le n,\ 0\le j\le i,\ 0\le a\le m,\ 0\le b\le t,
\]
nghĩa là để chọn \(j\) bài trong \(i\) bài đầu, số đĩa ít nhất cần dùng là
\(a\) đĩa cộng thêm \(b\) phút. Vì bài thứ \(i\) có hai khả năng, phát hành
hoặc không phát hành, ta có
\[
  g[i,j]=\min\{g[i-1,j],\ g[i-1,j-1]+\operatorname{long}[i]\}.
\]
Với \((a,b)+\operatorname{long}[i]=(a',b')\), cách tính là:
\[
\begin{array}{ll}
\operatorname{long}[i]\le t-b: & a'=a,\quad b'=b+\operatorname{long}[i],\\
\operatorname{long}[i]>t-b: & a'=a+1,\quad b'=\operatorname{long}[i].
\end{array}
\]
Điều kiện biên:
\[
  g[i,0]=(0,0) \qquad (0\le i\le n).
\]
Kết quả cần tìm là
\[
  \operatorname{ans}=\max\{k\mid g[n,k]\le (m-1,t)\}.
\]

Thuật toán cải tiến có \(O(n^2)\) trạng thái, mỗi trạng thái chuyển từ
\(O(1)\) trạng thái, mỗi lần chuyển mất \(O(1)\), nên tổng thời gian là
\(O(n^2)\). Đáng chú ý là độ phức tạp không gian cũng giảm từ \(O(mnt)\)
xuống \(O(n^2)\).

Qua tối ưu hóa bài này, ta thấy: các cách biểu diễn trạng thái khác nhau có
thể tạo ra những thuật toán quy hoạch động có hiệu năng rất khác nhau. Cải
tiến biểu diễn trạng thái có thể giảm tổng số trạng thái, từ đó hạ độ phức
tạp thời gian; đồng thời cũng có thể giảm độ phức tạp không gian. Vì vậy,
giảm tổng số trạng thái giữ vị trí quan trọng trong tối ưu hóa quy hoạch
động.

\subsubsection{Chọn hướng quy hoạch thích hợp}

Trong cài đặt quy hoạch động, có hai hướng quy hoạch chính: suy tiến và suy
lùi. Trong một số trường hợp, chọn hướng quy hoạch khác nhau cũng làm hiệu
suất thời gian của chương trình khác nhau. Nói chung, nếu trạng thái đầu xác
định còn trạng thái đích không xác định, nên xét dùng suy tiến; ngược lại,
nếu trạng thái đích xác định còn trạng thái đầu không xác định, nên xét dùng
suy lùi. Nếu cả trạng thái đầu và trạng thái đích đều đã xác định, thông
thường suy tiến và suy lùi đều dùng được; nhưng liệu có thể xét quy hoạch hai
chiều hay không?

Phương pháp tìm kiếm hai chiều đã rất quen thuộc. Tư tưởng chính là: khi
không gian trạng thái rất lớn, còn trạng thái đầu và đích đều xác định, lượng
trạng thái mở rộng tăng theo cấp số mũ. Để giảm quy mô trạng thái, ta mở rộng
từ cả trạng thái đầu lẫn trạng thái đích, rồi nhận được nghiệm tại nơi hai
phía gặp nhau.

Tư tưởng tối ưu trên có thể áp dụng cho quy hoạch động hay không? Hãy xét ví
dụ sau.

\paragraph{Ví dụ 2: Divide (Merc 2000).}

\paragraph{Mô tả bài toán.}
Có \(a[1..6]\) viên bi đá có giá trị lần lượt là \(1..6\). Cần chia chúng
thành hai phần sao cho tổng giá trị hai phần bằng nhau. Hỏi có thể làm được
không? Tổng số viên bi đá không vượt quá \(20000\).

\paragraph{Phân tích thuật toán.}
Đặt
\[
  S=\sum_{i=1}^{6} i\cdot a[i].
\]
Nếu \(S\) lẻ thì không thể chia; nếu không, đặt \(\operatorname{Mid}=S/2\),
bài toán trở thành: có thể chọn một số viên trong các viên đã cho sao cho
tổng giá trị bằng \(\operatorname{Mid}\) hay không.

Đây thực chất là bài toán hàm sinh, và dùng quy hoạch động là tương đương.
Đặt
\[
  m[i,j]\quad (0\le i\le6,\ 0\le j\le \operatorname{Mid})
\]
biểu thị có thể chọn một số viên có giá trị \(1..i\) sao cho tổng giá trị là
\(j\) hay không. Nếu được thì giá trị là \(\mathrm{true}\), ngược lại là
\(\mathrm{false}\). Công thức chuyển là
\[
  m[i,j]=m[i,j]\vee m[i-1,j-i\cdot k]\qquad (0\le k\le a[i]).
\]
Điều kiện biên là
\[
  m[i,0]=\mathrm{true}\qquad (0\le i\le6).
\]
Nếu \(m[i,\operatorname{Mid}]=\mathrm{true}\) với một \(0\le i\le6\) thì yêu
cầu bài toán thực hiện được; nếu không thì không thể.

Trong thuật toán trên, mỗi trạng thái có thể chuyển từ \(a[i]\) trạng thái,
mỗi lần chuyển mất \(O(1)\). Tổng số trạng thái là tổng số trạng thái có giá
trị \(\mathrm{true}\), thực chất chính là số hạng trong hàm sinh.

\paragraph{Tối ưu hóa thuật toán.}
Thực nghiệm cho thấy khi \(i\) nhỏ, do các loại giá trị viên bi còn đơn điệu
và số lượng ít, số trạng thái \(\mathrm{true}\) cũng ít. Nhưng khi \(i\) tăng,
số loại giá trị và số lượng viên bi tăng, số trạng thái \(\mathrm{true}\)
tăng nhanh, làm quá trình quy hoạch chậm lại và ảnh hưởng đến hiệu suất.

Mặt khác, ta chỉ quan tâm liệu có nhận được một trạng thái
\(\mathrm{true}\) có tổng giá trị \(\operatorname{Mid}\) hay không. Vì vậy có
thể quy hoạch từ hai hướng: lần lượt tính mọi tổng giá trị có thể thu được
khi chọn một số viên có giá trị \(1..3\), và mọi tổng giá trị có thể thu được
khi chọn một số viên có giá trị \(4..6\). Cuối cùng kiểm tra hai phía có tồn
tại hai tổng cộng lại bằng \(\operatorname{Mid}\) hay không; từ đó nhận được
nghiệm bài toán.

Công thức chuyển trạng thái cải tiến là
\[
\begin{array}{ll}
i\le3: & m[i,j]=m[i,j]\vee m[i-1,j-i\cdot k]\quad (1\le k\le a[i]),\\
i>3: & m[i,j]=m[i,j]\vee m[i+1,j-i\cdot k]\quad (1\le k\le a[i]).
\end{array}
\]
Điều kiện biên là
\[
  m[i,0]=\mathrm{true}\qquad (0\le i\le7).
\]
Nếu tồn tại \(k\) sao cho
\[
  m[3,k]=\mathrm{true},\qquad
  m[4,\operatorname{Mid}-k]=\mathrm{true},
\]
thì có thể thực hiện yêu cầu bài toán; nếu không thì không thể.

Trong PDF gốc, hình minh họa cho thấy sự khác biệt về tổng số trạng thái phải
tính giữa quy hoạch động hai chiều và quy hoạch động một chiều.

Nhìn lại quá trình tối ưu của bài này, có thể thấy bối cảnh thực tế của bài
rất giống bối cảnh của tìm kiếm hai chiều: cùng có không gian trạng thái lớn,
trạng thái đầu và trạng thái đích xác định, lượng trạng thái tăng nhanh, và
có thể kiểm tra giao nhau. Vì vậy, từ quá trình tối ưu bài này, ta thấy
phương pháp mở rộng hai chiều để giảm lượng trạng thái không chỉ thích hợp
với tìm kiếm, mà cũng thích hợp với quy hoạch động. Việc tìm thuộc tính chung
giữa những phương pháp giải bài khác nhau để mượn cùng một tư tưởng tối ưu có
thể giúp ta liên tục tạo ra phương pháp mới.

\subsection{Giảm số trạng thái chuyển của mỗi trạng thái}

Khi dùng quy hoạch động, việc tính trạng thái hiện tại đều dựa trên một số
quyết định và tham chiếu những trạng thái đã tính tương ứng. Quá trình đó gọi
là \term{chuyển trạng thái}. Vì vậy, số quyết định có thể đưa ra ở mỗi trạng
thái, cũng chính là số trạng thái mà mỗi trạng thái có thể chuyển từ đó, là
một yếu tố quan trọng quyết định độ phức tạp thời gian của thuật toán quy
hoạch động. Mục này thảo luận một số phương pháp giảm số trạng thái có thể
chuyển của mỗi trạng thái.

\subsubsection{Bất đẳng thức tứ giác và tính đơn điệu của quyết định}

\paragraph{Ví dụ 3: bài toán gộp đá (NOI 1995).}

\paragraph{Mô tả bài toán.}
Trên một sân chơi có một hàng \(n\) đống đá, \(n\le20\). Cần gộp các đống đá
theo thứ tự thành một đống. Mỗi lần chỉ được chọn hai đống kề nhau để gộp
thành một đống mới, và số viên đá của đống mới được tính là điểm của lần gộp
ấy.

Hãy lập chương trình tìm tổng điểm nhỏ nhất, tổng điểm lớn nhất và phương án
gộp tương ứng để gộp \(n\) đống đá thành một đống.

\paragraph{Phân tích thuật toán.}
Đây là một ứng dụng kinh điển của quy hoạch động. Vì bài toán điểm lớn nhất
và điểm nhỏ nhất tương tự nhau, ở đây chỉ thảo luận điểm nhỏ nhất. Đánh số
\(n\) đống đá theo thứ tự là \(1,2,\ldots,n\). Gọi số đá của các đống là
\(d[1..n]\). Biểu diễn trạng thái là
\[
  m[i,j]\quad (1\le i\le j\le n),
\]
biểu thị điểm nhỏ nhất nhận được khi gộp \(d[i..j]\). Công thức chuyển và
điều kiện biên là
\[
  m[i,j]=0\qquad (i=j),
\]
\[
  m[i,j]=\min_{i<k\le j}
  \left\{m[i,k-1]+m[k,j]+\sum_{l=i}^{j}d[l]\right\}
  \qquad (i<j).
\]
Đồng thời đặt \(s[i,j]=k\), biểu thị vị trí tách khi gộp, để sau khi tính ra
giá trị tối ưu có thể dựng lại nghiệm tối ưu.

Tổng \(\sum_{l=i}^{j}d[l]\) trong công thức có thể tính bằng tiền xử lý. Đặt
\[
  t[i]=\sum_{j=1}^{i}d[j]\quad (i=1..n),\qquad t[0]=0,
\]
khi đó
\[
  \sum_{l=i}^{j}d[l]=t[j]-t[i-1].
\]

Thuật toán trên có \(O(n^2)\) trạng thái, mỗi trạng thái chuyển từ \(O(n)\)
trạng thái, mỗi lần chuyển mất \(O(1)\), nên tổng thời gian là \(O(n^3)\).

\paragraph{Tối ưu hóa thuật toán.}
Khi hàm \(w[i,j]\) thỏa
\[
  w[i,j]+w[i',j']\le w[i',j]+w[i,j'],
  \qquad i\le i'\le j\le j',
\]
ta nói \(w\) thỏa \term{bất đẳng thức tứ giác}. Khi hàm \(w[i,j]\) thỏa
\[
  w[i',j]\le w[i,j'],\qquad i\le i'\le j\le j',
\]
ta nói \(w\) đơn điệu theo quan hệ bao hàm của đoạn.

Trong bài gộp đá, đặt
\[
  w[i,j]=\sum_{l=i}^{j}d[l].
\]
Khi đó \(w[i,j]\) thỏa bất đẳng thức tứ giác. Đồng thời, từ \(d[i]\ge0\) và
\(t[i]\ge0\), biết rằng \(w[i,j]\) cũng thỏa tính đơn điệu.

Công thức quy hoạch có thể viết lại thành
\[
  m[i,j]=0\qquad (i=j),
\]
\[
  m[i,j]=\min_{i<k\le j}\{m[i,k-1]+m[k,j]\}+w[i,j]\qquad (i<j). \tag{1}
\]

Với hàm đơn điệu \(w\) thỏa bất đẳng thức tứ giác, có thể suy ra hàm
\(m[i,j]\) do công thức truy hồi (1) định nghĩa cũng thỏa bất đẳng thức tứ
giác:
\[
  m[i,j]+m[i',j']\le m[i',j]+m[i,j'],
  \qquad i\le i'\le j\le j'.
\]
Tính chất này có thể chứng minh bằng quy nạp như sau.

Ta quy nạp theo ``độ dài'' \(l=j'-i\) trong bất đẳng thức tứ giác. Khi
\(i=i'\) hoặc \(j=j'\), bất đẳng thức hiển nhiên đúng. Vì vậy khi \(l\le1\),
hàm \(m\) thỏa bất đẳng thức tứ giác. Sau đó xét hai trường hợp.

\paragraph{Trường hợp 1: \(i<i'=j<j'\).}
Khi đó bất đẳng thức tứ giác trở thành bất đẳng thức tam giác ngược
\[
  m[i,j]+m[j,j']\le m[i,j'].
\]
Đặt
\[
  k=\max\{p\mid m[i,j']=m[i,p-1]+m[p,j']+w[i,j']\}.
\]
Lại chia thành hai khả năng \(k\le j\) và \(k>j\). Chỉ xét \(k\le j\); trường
hợp \(k>j\) tương tự. Khi đó
\[
\begin{aligned}
m[i,j]+m[j,j']
&\le w[i,j]+m[i,k-1]+m[k,j]+m[j,j']\\
&\le w[i,j']+m[i,k-1]+m[k,j]+m[j,j']\\
&\le w[i,j']+m[i,k-1]+m[k,j']\\
&=m[i,j'].
\end{aligned}
\]

\paragraph{Trường hợp 2: \(i<i'<j<j'\).}
Đặt
\[
\begin{aligned}
y&=\max\{p\mid m[i',j]=m[i',p-1]+m[p,j]+w[i',j]\},\\
z&=\max\{p\mid m[i,j']=m[i,p-1]+m[p,j']+w[i,j']\}.
\end{aligned}
\]
Lại cần xét hai khả năng \(z\le y\) và \(z>y\). Chỉ xét \(z\le y\); trường
hợp \(z>y\) tương tự. Từ \(i<z\le y\le j\), có
\[
\begin{aligned}
m[i,j]+m[i',j']
&\le w[i,j]+m[i,z-1]+m[z,j]
   +w[i',j']+m[i',y-1]+m[y,j']\\
&\le w[i,j']+w[i',j]+m[i',y-1]+m[i,z-1]+m[z,j]+m[y,j']\\
&\le w[i,j']+w[i',j]+m[i',y-1]+m[i,z-1]+m[y,j]+m[z,j']\\
&=m[i,j']+m[i',j].
\end{aligned}
\]
Từ đó \(m[i,j]\) thỏa bất đẳng thức tứ giác.

Đặt
\[
  s[i,j]=\max\{k\mid m[i,j]=m[i,k-1]+m[k,j]+w[i,j]\}.
\]
Từ việc \(m[i,j]\) thỏa bất đẳng thức tứ giác, có thể suy ra tính đơn điệu
của \(s[i,j]\):
\[
  s[i,j]\le s[i,j+1]\le s[i+1,j+1]\qquad (i\le j).
\]
Khi \(i=j\), tính đơn điệu hiển nhiên đúng, nên chỉ cần xét \(i<j\). Do tính
đối xứng, chỉ cần chứng minh \(s[i,j]\le s[i,j+1]\).

Đặt
\[
  m_k[i,j]=m[i,k-1]+m[k,j]+w[i,j].
\]
Để chứng minh \(s[i,j]\le s[i,j+1]\), chỉ cần chứng minh với mọi
\(i<k\le k'\le j\), nếu \(m_{k'}[i,j]\le m_k[i,j]\), thì
\[
  m_{k'}[i,j+1]\le m_k[i,j+1].
\]
Thực ra ta có thể chứng minh bất đẳng thức mạnh hơn:
\[
  m_k[i,j]-m_{k'}[i,j]\le m_k[i,j+1]-m_{k'}[i,j+1],
\]
tức là
\[
  m_k[i,j]+m_{k'}[i,j+1]\le m_k[i,j+1]+m_{k'}[i,j].
\]
Khai triển theo định nghĩa truy hồi và sắp xếp lại, ta được
\[
  m[k,j]+m[k',j+1]\le m[k',j]+m[k,j+1],
\]
đúng là bất đẳng thức tứ giác khi \(k\le k'\le j<j+1\). Do đó, khi \(w\) thỏa
bất đẳng thức tứ giác, hàm \(s[i,j]\) có tính đơn điệu.

Vì vậy, sử dụng tính đơn điệu của \(s[i,j]\), ta có công thức chuyển trạng
thái tối ưu:
\[
  m[i,j]=0\qquad (i=j),
\]
\[
  m[i,j]=
  \min_{s[i,j-1]\le k\le s[i+1,j]}
  \{m[i,k-1]+m[k,j]\}+w[i,j]\qquad (i<j).
\]
Bằng phương pháp tương tự có thể chứng minh bài toán điểm lớn nhất cũng dùng
được cùng cách tối ưu.

Thời gian tính của công thức chuyển cải tiến là
\[
\begin{aligned}
&O\left(\sum_{i=1}^{n-1}\sum_{j=i+1}^{n}
  (1+s[i+1,j]-s[i,j-1])\right)\\
&\quad =
O\left(\sum_{i=1}^{n-1}
  (n-i+s[i+1,n]-s[1,n-i])\right)
=O(n^2).
\end{aligned}
\]

Phương pháp trên dùng bất đẳng thức tứ giác để suy ra tính đơn điệu của quyết
định tối ưu, từ đó giảm số trạng thái chuyển của mỗi trạng thái và hạ độ phức
tạp thời gian. Phương pháp này có tính phổ quát. Với những bài quy hoạch động
có công thức chuyển tương tự (1), và \(w[i,j]\) thỏa bất đẳng thức tứ giác,
đều có thể dùng cùng cách tối ưu, chẳng hạn cây nhị phân tìm kiếm tối ưu
(NOI 1996). Dưới đây là một ví dụ khác.

\paragraph{Ví dụ 4: Post Office (IOI 2000).}

\paragraph{Mô tả bài toán.}
Cho \(n\) ngôi làng trên một đường thẳng, có tọa độ đôi một khác nhau và được
đưa ra theo thứ tự tăng dần. Cần chọn \(p\) làng để xây bưu cục. Mỗi làng sử
dụng bưu cục gần nó nhất, sao cho tổng khoảng cách từ mọi làng đến bưu cục
mà chúng sử dụng là nhỏ nhất.

Hãy lập chương trình tính tổng khoảng cách nhỏ nhất và phương án xây bưu cục.

\paragraph{Phân tích thuật toán.}
Đây cũng là một bài quy hoạch động. Đánh số \(n\) làng theo thứ tự tọa độ
tăng dần là \(1,2,\ldots,n\), tọa độ là \(d[1..n]\). Biểu diễn trạng thái là:
\[
  m[i,j]=\text{tổng khoảng cách nhỏ nhất khi xây } i
  \text{ bưu cục cho } j \text{ làng đầu}.
\]
Vì vậy \(m[p,n]\) là nghiệm bài toán. Công thức chuyển và điều kiện biên là
\[
  m[1,j]=w[1,j],
\]
\[
  m[i,j]=\min_{i-1\le k\le j-1}
  \{m[i-1,k]+w[k+1,j]\}\qquad (i\le j).
\]
Trong đó \(w[i,j]\) là tổng khoảng cách nhỏ nhất khi xây một bưu cục trong
các làng \(d[i..j]\). Có thể chứng minh rằng khi chỉ xây một bưu cục, nghiệm
tối ưu nằm ở trung vị, tức nếu làng xây bưu cục là \(k\) thì
\[
  k=\left\lfloor\frac{i+j}{2}\right\rfloor
  \quad\text{hoặc}\quad
  k=\left\lceil\frac{i+j}{2}\right\rceil.
\]
Do đó
\[
  w[i,j]=\sum_{l=i}^{j}|d[l]-d[k]|,\qquad
  k=\left\lfloor\frac{i+j}{2}\right\rfloor
  \text{ hoặc }
  k=\left\lceil\frac{i+j}{2}\right\rceil.
\]
Đồng thời đặt \(s[i,j]=k\), ghi lại số làng do \(i-1\) bưu cục trước phục vụ,
để sau khi tính tổng khoảng cách nhỏ nhất có thể dựng lại phương án tối ưu.

Trong thuật toán trên, \(w[i,j]\) có thể được tiền xử lý trong \(O(n)\) thời
gian để sau đó tính trong \(O(1)\). Vì vậy thuật toán có \(O(np)\) trạng
thái, mỗi trạng thái chuyển từ \(O(n)\) trạng thái, mỗi lần chuyển mất
\(O(1)\), và tổng thời gian là \(O(pn^2)\).

\paragraph{Tối ưu hóa thuật toán.}
Công thức chuyển trạng thái của bài này rất giống (1), nên ta phỏng đoán
quyết định cũng thỏa tính đơn điệu:
\[
  s[i-1,j]\le s[i,j]\le s[i,j+1].
\]
Trước hết chứng minh \(w\) thỏa bất đẳng thức tứ giác:
\[
  w[i,j]+w[i',j']\le w[i',j]+w[i,j'],
  \qquad i\le i'\le j\le j'.
\]
Đặt
\[
  y=\left\lfloor\frac{i'+j}{2}\right\rfloor,\qquad
  z=\left\lfloor\frac{i+j'}{2}\right\rfloor.
\]
Chia thành hai trường hợp \(z\le y\) và \(z>y\); ở đây chỉ xét \(z\le y\),
trường hợp còn lại tương tự. Từ \(i\le z\le y\le j\), có
\[
\begin{aligned}
w[i,j]+w[i',j']
&\le \sum_{l=i}^{j}|d[l]-d[z]|
   +\sum_{l=i'}^{j'}|d[l]-d[y]|\\
&\le \sum_{l=i}^{j}|d[l]-d[z]|
   +\sum_{l=i'}^{j'}|d[l]-d[y]|\\
&\quad +\sum_{l=j+1}^{j'}|d[l]-d[z]|
   -\sum_{l=j+1}^{j'}|d[l]-d[y]|\\
&=\sum_{l=i}^{j'}|d[l]-d[z]|
  +\sum_{l=i'}^{j}|d[l]-d[y]|\\
&=w[i',j]+w[i,j'].
\end{aligned}
\]

Tiếp theo, dùng quy nạp toán học để chứng minh hàm \(m\) cũng thỏa bất đẳng
thức tứ giác. Quy nạp theo ``độ dài'' \(l=j'-i\). Khi \(i=i'\) hoặc \(j=j'\),
bất đẳng thức hiển nhiên đúng, nên khi \(l\le1\), \(m\) thỏa bất đẳng thức tứ
giác. Sau đó xét hai trường hợp.

\paragraph{Trường hợp 1: \(i<i'=j<j'\).}
Cần chứng minh
\[
  m[i,j]+m[j,j']\le m[i,j'].
\]
Đặt
\[
  k=\max\{p\mid m[i,j']=m[i,p-1]+m[p,j']+w[i,j']\}.
\]
Lại chia thành \(k\le j\) và \(k>j\). Chỉ xét \(k\le j\); trường hợp còn lại
tương tự:
\[
\begin{aligned}
m[i,j]+m[j,j']
&\le m[i-1,k]+w[k+1,j]+m[j-1,j-1]+w[j,j']\\
&\le m[i-1,k]+w[k+1,j']\\
&=m[i,j'].
\end{aligned}
\]

\paragraph{Trường hợp 2: \(i<i'<j<j'\).}
Đặt
\[
\begin{aligned}
y&=\max\{p\mid m[i',j]=m[i'-1,p]+w[p+1,j]\},\\
z&=\max\{p\mid m[i,j']=m[i-1,p]+w[p+1,j']\}.
\end{aligned}
\]
Lại chia thành \(z\le y\) và \(z>y\).

Khi \(z\le y<j<j'\):
\[
\begin{aligned}
m[i,j]+m[i',j']
&\le m[i'-1,y]+w[y+1,j']+m[i-1,z]+w[z+1,j]\\
&\le m[i'-1,y]+m[i-1,z]+w[y+1,j]+w[z+1,j']\\
&=m[i',j]+m[i,j'].
\end{aligned}
\]
Khi \(i-1<i'-1\le y<z<j'\):
\[
\begin{aligned}
m[i,j]+m[i',j']
&\le m[i-1,y]+w[y+1,j]+m[i'-1,z]+w[z+1,j']\\
&\le m[i-1,z]+m[i'-1,y]+w[y+1,j]+w[z+1,j']\\
&=m[i,j']+m[i',j].
\end{aligned}
\]

Cuối cùng chứng minh quyết định \(s[i,j]\) thỏa tính đơn điệu. Để tiện thảo
luận, đặt
\[
  m_k[i,j]=m[i-1,k]+w[k+1,j].
\]
Trước hết chứng minh \(s[i-1,j]\le s[i,j]\). Chỉ cần chứng minh: với mọi
\(i\le k<k'<j\), nếu \(m_{k'}[i-1,j]\le m_k[i-1,j]\), thì
\[
  m_{k'}[i,j]\le m_k[i,j].
\]
Tương tự, ta chứng minh bất đẳng thức mạnh hơn:
\[
  m_k[i-1,j]-m_{k'}[i-1,j]\le m_k[i,j]-m_{k'}[i,j],
\]
tức là
\[
  m_k[i-1,j]+m_{k'}[i,j]\le m_k[i,j]+m_{k'}[i-1,j].
\]
Khai triển theo định nghĩa truy hồi và sắp xếp lại được
\[
  m[i-2,k]+m[i-1,k']\le m[i-1,k]+m[i-2,k'],
\]
chính là bất đẳng thức tứ giác của \(m\) khi \(i-2<i-1<k<k'\).

Tiếp đó chứng minh \(s[i,j]\le s[i,j+1]\). Tương tự phần trên, với
\(k<k'<j\), chỉ cần chứng minh bất đẳng thức mạnh hơn:
\[
  m_k[i,j]-m_{k'}[i,j]\le m_k[i,j+1]-m_{k'}[i,j+1],
\]
tức là
\[
  m_k[i,j]+m_{k'}[i,j+1]\le m_k[i,j+1]+m_{k'}[i,j].
\]
Khai triển theo định nghĩa truy hồi và sắp xếp lại được
\[
  w[k+1,j]+w[k'+1,j+1]\le w[k+1,j+1]+w[k'+1,j],
\]
chính là bất đẳng thức tứ giác của \(w\) khi \(k+1<k'+1<j<j+1\).

Từ đó, quyết định \(s[i,j]\) của bài toán có tính đơn điệu. Công thức chuyển
tối ưu là
\[
  m[1,j]=w[1,j],
\]
\[
  m[i,j]=
  \min_{s[i-1,j]\le k\le s[i,j+1]}
  \{m[i-1,k]+w[k+1,j]\}\qquad (i\le j),
\]
và \(s[i,j]=k\). Như đã nêu ở trên, thuật toán tối ưu có thời gian \(O(np)\).

Bản chất của tối ưu bằng bất đẳng thức tứ giác là tận dụng đầy đủ kết quả.
Thông qua phân tích quan hệ đặc biệt giữa các giá trị trạng thái, nó suy ra
tính đơn điệu của quyết định tối ưu. Nhờ đó, khi tính trạng thái hiện tại, ta
tận dụng những quyết định tối ưu đã đưa ra ở các trạng thái đã tính, giảm
lượng quyết định hiện tại. Điều này gợi ý rằng khi dùng quy hoạch động, ta
không chỉ có thể tận dụng đầy đủ giá trị trạng thái, mà còn có thể tận dụng
đầy đủ quyết định tối ưu. Thực chất đây là một góc nhìn khác của việc giảm
dư thừa.

\subsubsection{Tối ưu hóa lượng quyết định}

Thông qua phân tích các tính chất mà nghiệm tối ưu của bài toán phải có, từ
đó thu hẹp tập quyết định có khả năng sinh ra nghiệm tối ưu, ta cũng giảm
được số trạng thái có thể chuyển của mỗi trạng thái.

Bài thêm dấu trong NOI 1996 mà mọi người quen thuộc chính là xuất phát từ
nguyên tắc ``tổng thu được là nhỏ nhất'', chỉ thêm dấu gần các điểm chia đều,
từ đó giảm mạnh số trạng thái chuyển của mỗi trạng thái và hạ độ phức tạp
thời gian. Hãy xét thêm một ví dụ.

\paragraph{Ví dụ 5: điểm lớn nhất trong bài gộp đá.}

\paragraph{Mô tả bài toán.}
Xem Ví dụ 3; trong ví dụ này chỉ xét bài toán điểm lớn nhất.

\paragraph{Phân tích thuật toán.}
Đánh số \(n\) đống đá theo thứ tự là \(1,2,\ldots,n\), số đá của các đống là
\(d[1..n]\). Biểu diễn trạng thái là
\[
  m[i,j]\quad (1\le i\le j\le n),
\]
biểu thị điểm lớn nhất nhận được khi gộp \(d[i..j]\). Công thức chuyển và
điều kiện biên là
\[
  m[i,j]=0\qquad (i=j),
\]
\[
  m[i,j]=\max_{i<k\le j}
  \left\{m[i,k-1]+m[k,j]+\sum_{l=i}^{j}d[l]\right\}
  \qquad (i<j).
\]
Đồng thời đặt \(s[i,j]=k\), biểu thị vị trí tách khi gộp, để sau khi tính ra
giá trị tối ưu có thể dựng lại nghiệm. Thuật toán có thời gian \(O(n^3)\).

\paragraph{Tối ưu hóa thuật toán.}
Phân tích kỹ bài toán có thể thấy \(s[i,j]\) hoặc bằng \(i+1\), hoặc bằng
\(j\), tức là
\[
  \max_{i<k\le j}\{m[i,k-1]+m[k,j]\}
  =\max\{m[i,j-1],m[i+1,j]\}\qquad (i<j).
\]

Có thể chứng minh bằng phản chứng. Giả sử vị trí tách làm \(m[i,j]\) đạt lớn
nhất là \(p\), với \(i+1<p<j\). Đặt
\[
  y=\sum_{l=i}^{p-1}a[l],\qquad
  z=\sum_{l=p}^{j}a[l].
\]
Chia thành hai trường hợp.

Trường hợp 1: \(y\ge z\). Vì \(p<j\), đặt \(s[p,j]=k\). Khi đó cách gộp tương
ứng có thể viết là
\[
  ((a[i]\ldots a[p-1])((a[p]\ldots a[k-1])(a[k]\ldots a[j]))).
\]
Điểm tương ứng là
\[
  T=m[i,p-1]+m[p,k-1]+m[k,j]+y+z+z. \tag{2}
\]
Xét một phương án gộp khác với \(s'[i,j]=k\) và \(s'[i,k]=p\):
\[
  (((a[i]\ldots a[p-1])(a[p]\ldots a[k-1]))(a[k]\ldots a[j])).
\]
Điểm tương ứng là
\[
  T'=m[i,p-1]+m[p,k-1]+m[k,j]+y+y+z+\sum_{l=p}^{k-1}a[l]. \tag{3}
\]
Từ \(y\ge z\) suy ra \(T<T'\), mâu thuẫn với giả thiết rằng vị trí tách làm
\(m[i,j]\) đạt lớn nhất là \(p\). Trường hợp 2, \(y<z\), chứng minh tương tự.

Vì vậy công thức chuyển trạng thái được tối ưu thành
\[
  m[i,j]=0\qquad (i=j),
\]
\[
  m[i,j]=\max\{m[i,j-1],m[i+1,j]\}+\sum_{l=i}^{j}d[l]\qquad (i<j).
\]
Sau tối ưu, số trạng thái chuyển của mỗi trạng thái giảm xuống \(O(1)\), tổng
thời gian cũng giảm xuống \(O(n^2)\).

Quá trình tối ưu bài này dựa vào phân tích tính chất của nghiệm tối ưu để tìm
điều kiện cần mà quyết định tối ưu phải thỏa. Điều này rất giống tư tưởng cắt
tỉa tối ưu trong tìm kiếm. Từ đó ta lại thấy cùng một tư tưởng tối ưu có thể
được áp dụng cho những phương pháp thiết kế thuật toán khác nhau. Đồng thời,
ta cũng thấy tối ưu hóa quy hoạch động phải xây dựng trên cơ sở phân tích bài
toán toàn diện và tỉ mỉ. Chỉ khi phân tích sâu thuộc tính của bài toán và đào
sâu bản chất của nó, ta mới tối ưu được thuật toán.

\subsubsection{Tổ chức trạng thái hợp lý}

Trong quá trình giải bằng quy hoạch động, ta cần liên tục tham chiếu những
trạng thái đã tính. Vì vậy, tổ chức hợp lý các trạng thái đã tính có lợi cho
việc nâng cao hiệu suất thời gian của quy hoạch động.

\paragraph{Ví dụ 6: tìm dãy con tăng dài nhất.}

\paragraph{Mô tả bài toán.}
Cho một dãy gồm \(n\) số \(x[1..n]\). Hãy tìm dãy con tăng dài nhất của nó,
tức là tìm \(m\) lớn nhất và các chỉ số \(a_1,a_2,\ldots,a_m\) sao cho
\[
  a_1<a_2<\cdots<a_m
  \quad\text{và}\quad
  x[a_1]<x[a_2]<\cdots<x[a_m].
\]

\paragraph{Phân tích thuật toán.}
Đây cũng là một ứng dụng kinh điển của quy hoạch động. Biểu diễn trạng thái
là
\[
  m[i]\quad (1\le i\le n),
\]
biểu thị độ dài dãy con tăng dài nhất kết thúc tại \(x[i]\). Nghiệm bài toán
là \(\max\{m[i]\mid 1\le i\le n\}\). Công thức chuyển và điều kiện biên là
\[
  m[i]=1+\max\{0,m[k]\mid x[k]<x[i],\ 1\le k<i\}.
\]
Đồng thời khi \(m[i]>1\), đặt \(p[i]=k\) để ghi quyết định tối ưu, nhằm dựng
lại dãy con tăng dài nhất sau khi tính được giá trị tối ưu.

Thuật toán trên có \(O(n)\) trạng thái, mỗi trạng thái chuyển từ nhiều nhất
\(O(n)\) trạng thái, mỗi lần chuyển mất \(O(1)\), nên tổng thời gian là
\(O(n^2)\).

\paragraph{Tối ưu hóa thuật toán.}
Trước hết xét hai tình huống.

Thứ nhất, nếu \(x[i]<x[j]\) và \(m[i]=m[j]\), thì không cần giữ trạng thái
\(m[j]\). Vì mọi trạng thái \(m[k]\) \((k>j,\ k>i)\) có thể chuyển từ
\(m[j]\) đều thỏa \(x[k]>x[j]>x[i]\), nên cũng có thể chuyển từ \(m[i]\). Mặt
khác, có những trạng thái \(m[k]\) có thể chuyển từ \(m[i]\), với
\(x[j]>x[k]>x[i]\), lại không thể chuyển từ \(m[j]\). Vì vậy, trong mọi trạng
thái có cùng giá trị, chỉ cần giữ trạng thái có giá trị phần tử cuối nhỏ
nhất.

Thứ hai, nếu \(x[i]<x[j]\) và \(m[i]>m[j]\), thì không cần giữ trạng thái
\(m[j]\). Vì mọi trạng thái \(m[k]\) \((k>j,\ k>i)\) có thể chuyển từ
\(m[j]\) đều thỏa \(x[k]>x[j]>x[i]\), nên cũng có thể chuyển từ \(m[i]\). Hơn
nữa \(m[i]>m[j]\), nên \(m[k]\ge m[i]+1>m[j]+1\); do đó chuyển từ \(m[j]\)
không có ý nghĩa.

Tổng hợp hai điểm trên, ta nhận được điều kiện cần để trạng thái \(m[k]\) cần
được giữ lại: không tồn tại \(i\) sao cho \(x[i]<x[k]\) và \(m[i]\ge m[k]\).
Vì vậy, trong các trạng thái giữ lại không có hai trạng thái cùng giá trị; và
khi giá trị trạng thái tăng, giá trị phần tử cuối cũng tăng đơn điệu.

Nói cách khác, giả sử tập trạng thái hiện đang giữ là \(S\), thì \(S\) có
tính chất \(D\): với mọi \(i\in S\), \(j\in S\), \(i\ne j\), ta có
\(m[i]\ne m[j]\); nếu \(m[i]<m[j]\) thì \(x[i]<x[j]\), ngược lại thì
\(x[i]>x[j]\).

Bây giờ xét chuyển trạng thái. Giả sử đã tính \(m[1..i-1]\), tập trạng thái
đang giữ là \(S\), và cần tính \(m[i]\).
\begin{enumerate}
  \item Nếu tồn tại trạng thái \(k\in S\) sao cho \(x[k]=x[i]\), thì trạng
  thái \(m[i]\) chắc chắn không cần giữ, cũng không cần tính. Thật vậy, giả sử
  \(m[i]=m[j]+1\), thì \(x[j]<x[i]=x[k]\), \(j\in S\), \(j\ne k\). Do đó
  \(m[j]<m[k]\), suy ra \(m[i]=m[j]+1\le m[k]\), nên trạng thái \(m[i]\)
  không cần giữ.
  \item Nếu không, \(m[i]=1+\max\{m[j]\mid x[j]<x[i],\ j\in S\}\). Các chỉ số
  \(j\) thỏa điều kiện cũng thỏa
  \[
    x[j]=\max\{x[k]\mid x[k]<x[i],\ k\in S\}.
  \]
  Đồng thời thêm trạng thái \(i\) vào \(S\).
  \item Nếu trường hợp 2 xảy ra, ta đã thêm một trạng thái vào \(S\). Để giữ
  tính chất của \(S\), cần bảo trì \(S\). Nếu tồn tại trạng thái \(k\in S\)
  sao cho \(m[i]=m[k]\), thì ta có \(x[i]<x[k]\) và
  \[
    x[k]=\min\{x[j]\mid x[j]>x[i],\ j\in S\}.
  \]
  Khi đó phải xóa trạng thái \(k\) khỏi \(S\).
\end{enumerate}

Từ đó có thuật toán cải tiến:
\begin{verbatim}
For i:=1 to n do
{
  Tim phan tu k co gia tri x lon nhat khong vuot qua x[i] trong S;
  if x[k]<x[i] then
  {
     m[i]:=m[k]+1;
     Chen trang thai i vao S;
     Tim phan tu j co gia tri x nho nhat lon hon x[i] trong S;
     if m[j]=m[i] then xoa trang thai j khoi S;
  }
}
\end{verbatim}

Từ tính chất \(D\) và mô tả thuật toán, có thể thấy \(S\) thực chất là một tập
có thứ tự theo khóa \(x\) (và cũng theo khóa \(m\)). Nếu dùng cây cân bằng để
cài đặt tập có thứ tự \(S\), thuật toán có độ phức tạp \(O(n\log_2 n)\). Sau
tối ưu, số trạng thái chuyển của mỗi trạng thái chỉ là \(O(1)\), còn thời
gian mỗi lần chuyển trở thành \(O(\log_2 n)\). Điều này cũng thể hiện mâu
thuẫn giữa các yếu tố tối ưu đã nói ở trên; nhìn toàn cục, độ phức tạp thời
gian của thuật toán vẫn giảm.

Nhìn lại quá trình tối ưu bài này, trước hết ta phân tích quan hệ giữa các
trạng thái để giảm số trạng thái cần giữ, đồng thời phát hiện tính đơn điệu
của những trạng thái cần giữ. Từ đó giảm số trạng thái có thể chuyển của mỗi
trạng thái, rồi dùng cấu trúc dữ liệu hiệu quả là cây cân bằng để tổ chức
những trạng thái hiện đang giữ, hoàn thành tối ưu thuật toán.

Qua bài này, ta thấy giảm số trạng thái được giữ lại và tổ chức hợp lý các
trạng thái đã tính có thể giúp giảm số trạng thái có thể chuyển của mỗi trạng
thái. Đồng thời, chọn cấu trúc dữ liệu thích hợp cũng là một nguyên tắc quan
trọng của tối ưu thuật toán. Ở phần sau, ta sẽ còn thấy cách dùng cấu trúc dữ
liệu để tối ưu thuật toán.

\subsubsection{Làm mịn chuyển trạng thái}

``Làm mịn chuyển trạng thái'' nghĩa là chia một lần chuyển trạng thái ban đầu
thành nhiều lần chuyển trạng thái nhỏ hơn, nhằm giảm tổng số lần chuyển trạng
thái. Trước khi tối ưu, quyết định của bài toán thường là quyết định phức hợp,
tức là một cách sắp xếp của nhiều quyết định con. Vì vậy quy mô quyết định
lớn, và số trạng thái có thể chuyển của mỗi trạng thái cũng nhiều. Cách tối
ưu là chia mỗi quyết định phức hợp thành nhiều quyết định con, rồi thêm một
trạng thái sau mỗi quyết định con. Nhờ vậy, quyết định con phía sau chỉ tiếp
tục chuyển khi các quyết định con phía trước đã đạt nghiệm tối ưu. Sau tối
ưu, tuy tổng số trạng thái tăng, nhưng tổng số lần chuyển trạng thái giảm, nên
tổng độ phức tạp của thuật toán giảm.

Cần chú ý rằng tối ưu nói trên phải thỏa một điều kiện: các quyết định con
của mỗi quyết định phức hợp ban đầu cũng phải thỏa nguyên lý tối ưu và tính
không hậu hiệu. Nói cách khác, các quyết định con của một quyết định phức hợp
tối ưu cũng là quyết định tối ưu; các quyết định con phía trước không ảnh
hưởng đến các quyết định con phía sau.

Phương pháp tối ưu trên một lần nữa thể hiện rằng để tối ưu một yếu tố có thể
phải trả giá bằng việc tăng một yếu tố khác. Tuy nhiên, hạ tổng độ phức tạp
thời gian của thuật toán mới là mục tiêu thật sự.

\subsection{Giảm thời gian chuyển trạng thái}

Ta biết chuyển trạng thái là thao tác cơ bản của quy hoạch động. Vì vậy, giảm
thời gian cần cho mỗi lần chuyển trạng thái có ý nghĩa quan trọng đối với
việc nâng cao hiệu suất thời gian của thuật toán.

Chuyển trạng thái chủ yếu gồm hai phần:
\begin{enumerate}
  \item Đưa ra quyết định: dùng trạng thái hiện tại và quyết định được chọn
  để tính ra trạng thái cần tham chiếu.
  \item Tính công thức truy hồi: tính giá trị trạng thái hiện tại theo công
  thức truy hồi, trong đó thao tác chính là tính các hạng hằng.
\end{enumerate}
Mục này lần lượt thảo luận một số phương pháp nâng hiệu suất thời gian của
hai phần ấy.

\subsubsection{Giảm thời gian quyết định}

\paragraph{Ví dụ 7: LOSTCITY (NOI 2000).}

\paragraph{Mô tả bài toán.}
Cho một bảng từ, một số quy tắc ngữ pháp xác định và một bài văn:
\begin{itemize}
  \item Bài văn và bảng từ chỉ chứa 26 chữ cái thường tiếng Anh \(a\ldots z\).
  \item Từ trong bảng từ chỉ có ba từ loại: danh từ, động từ và phó từ. Các
  từ cùng từ loại đôi một khác nhau. Độ dài mỗi từ không vượt quá 20.
  \item Quy tắc ngữ pháp có thể mô tả ngắn gọn như sau: cụm danh từ gồm một
  số lượng tùy ý phó từ tiền tố rồi đến một danh từ; cụm động từ gồm một số
  lượng tùy ý phó từ tiền tố rồi đến một động từ; câu bắt đầu bằng một cụm
  danh từ, sau đó các cụm danh từ và cụm động từ nối xen kẽ nhau.
  \item Độ dài bài văn không vượt quá 1000. Biết rằng bài văn do hữu hạn câu
  tạo thành, và mỗi câu chỉ chứa hữu hạn từ.
\end{itemize}
Hãy lập chương trình chia bài văn này thành ít câu nhất; trong điều kiện đó,
số từ được chia ra cũng phải ít nhất.

\paragraph{Phân tích thuật toán.}
Đây cũng là một bài quy hoạch động. Lần lượt dùng \(v,u,a\) để biểu thị động
từ, danh từ và phó từ; bài văn cho trước là \(L[1..M]\). Biểu diễn trạng thái
là:
\[
\begin{array}{ll}
F(v,i): & \text{phương án tối ưu khi chia } i \text{ ký tự đầu của } L\\
        & \text{thành phần kết thúc bằng động từ; khi } i\ne M\\
        & \text{có thể kèm một số phó từ hậu tố tùy ý,}\\[0.2em]
F(u,i): & \text{phương án tối ưu khi chia } i \text{ ký tự đầu của } L\\
        & \text{thành phần kết thúc bằng danh từ; khi } i\ne M\\
        & \text{có thể kèm một số phó từ hậu tố tùy ý.}
\end{array}
\]
Mỗi trạng thái ghi cặp gồm số câu và số từ. Phương án phân tách trước đó chỉ
ảnh hưởng đến quyết định sau này thông qua từ loại của từ không phải phó từ
cuối cùng, nên biểu diễn trạng thái này thỏa tính không hậu hiệu.

Công thức chuyển trạng thái là
\[
\begin{aligned}
F(v,i)=\min\{&
  F(u,j)+(0,1), && L(j+1..i)\text{ là động từ};\\
& F(v,j)+(0,1), && L(j+1..i)\text{ là phó từ},\ i\ne M\},
\end{aligned}
\]
\[
\begin{aligned}
F(u,i)=\min\{&
  F(u,j)+(1,1), && L(j+1..i)\text{ là danh từ};\\
& F(v,j)+(0,1), && L(j+1..i)\text{ là danh từ};\\
& F(u,j)+(0,1), && L(j+1..i)\text{ là phó từ},\ i\ne M\}.
\end{aligned}
\]
Điều kiện biên:
\[
  F(v,0)=(1,0),\qquad F(u,0)=(\infty,\infty).
\]
Nghiệm bài toán là
\[
  \min\{F(v,M),F(u,M)\}.
\]

Trong thuật toán trên, tổng số trạng thái là \(O(M)\), mỗi trạng thái chuyển
từ nhiều nhất 20 trạng thái. Khi chuyển trạng thái, cần tra từ loại của
\(L[j+1..i]\), rồi dựa vào từ loại để đưa ra quyết định tương ứng và tham
chiếu trạng thái tương ứng. Dưới đây so sánh độ phức tạp thời gian khi dùng
các phương pháp khác nhau để tra từ loại của \(L[j+1..i]\).

\paragraph{Cài đặt thuật toán.}
Giả sử quy mô bảng từ là \(N\). Trước hết tiền xử lý bảng từ: sắp các từ theo
thứ tự từ điển và gộp những từ có nhiều từ loại. Khi tra từ loại, có thể dùng
các phương pháp sau.

\begin{enumerate}
  \item \term{Tìm tuần tự.} Trong trường hợp xấu nhất cần duyệt toàn bộ bảng
  từ, nên thời gian xấu nhất là \(O(20NM)\). Số phép so sánh nhiều nhất có
  thể đạt \(1000\cdot 5000\cdot 20=10^8\), hiệu suất thấp khi dữ liệu lớn.
  \item \term{Tìm nhị phân.} Thời gian xấu nhất là \(O(20M\log_2 N)\); số
  phép so sánh nhiều nhất giảm xuống khoảng \(5000\cdot20\cdot\log_2 1000
  \approx 10^6\), hoàn toàn chấp nhận được.
  \item \term{Bảng băm.} Đây là phương pháp hiệu quả nhất thường dùng cho
  tra cứu tập hợp. Trước hết cứ bốn ký tự thì gấp lại cộng để tính khóa
  \(k\), sau đó dùng băm kép để tính giá trị hàm băm \(h(k)\). Với phương
  pháp này, sau \(O(N)\) thời gian tiền xử lý để xây bảng băm, mỗi lần tra
  chỉ cần \(O(1)\), nên thời gian thuật toán là \(O(20M+N)=O(M)\).
  \item \term{Cây tra cứu.} Với tập hợp có phần tử là chuỗi, còn có một cách
  hiệu quả hơn là dùng cây tra cứu. Vì tất cả từ cần tra khi chuyển trạng
  thái tại một trạng thái đều nằm trên cùng một đường từ gốc đến lá của cây,
  nếu chọn thao tác đi từ gốc theo một đường đến lá làm thao tác cơ bản, thì
  nhiều nhất 20 lần tra từ của mỗi trạng thái chỉ cần \(O(1)\). Ngoài ra,
  xây cây tra cứu mất \(O(N)\). Vì vậy tổng thời gian vẫn là \(O(M)\), nhưng
  do hằng số nhỏ hơn phương pháp dùng bảng băm, tốc độ đo thực tế là nhanh
  nhất.
\end{enumerate}

Từ quá trình tối ưu bài này có thể thấy: dùng cấu trúc dữ liệu đúng là một
nguyên tắc quan trọng của tối ưu thuật toán, và điều này cũng đúng trong tối
ưu thuật toán quy hoạch động. Phương pháp 3 dùng bảng băm, một cấu trúc dữ
liệu tra cứu tập hợp hiệu quả; phương pháp 4 dùng cây tra cứu chuyên biệt
hơn, nhờ đó nâng cao hiệu suất thời gian của thuật toán.

\subsubsection{Giảm thời gian tính công thức truy hồi}

Thao tác chính khi tính công thức truy hồi là tính các hạng hằng. Vì vậy,
giảm thời gian cần để tính công thức truy hồi chủ yếu là giảm thời gian tính
các hạng hằng.

\paragraph{Ví dụ 8: tuần tra đường cao tốc (CTSC 2000).}

\paragraph{Mô tả bài toán.}
Trên một đường cao tốc không có ngã rẽ có \(n\) trạm gác, \(n\le50\). Khoảng
cách giữa hai trạm gác kề nhau đều là \(10\) km. Tốc độ thấp nhất của mọi xe
trên đường là \(60\) km/h, tốc độ cao nhất là \(120\) km/h, và xe chỉ được đổi
tốc độ tại trạm gác.

Có \(m\) xe tuần tra, \(m\le300\). Xe tuần tra thứ \(r\) xuất phát tại thời
điểm \(T_r\) từ trạm gác thứ \(n_r\), chạy đều đến trạm gác thứ \(n_r+1\),
thời gian trên đường là \(t_r\) giây.

Hai xe gặp nhau nếu giữa chúng xảy ra hiện tượng vượt xe hoặc chúng đồng thời
đến một trạm gác. Hãy tìm số xe tuần tra ít nhất mà một xe xuất phát đúng
6:00 từ trạm gác thứ 1 đi đến trạm gác thứ \(n\) sẽ gặp, và trong trường hợp
ấy thời điểm đến trạm gác thứ \(n\) sớm nhất. Giả sử thời điểm mọi xe đến
trạm gác đều là số nguyên giây.

\paragraph{Phân tích thuật toán.}
Bài này cũng giải bằng quy hoạch động. Biểu diễn trạng thái là:
\[
  F(i,T)=\text{số lần gặp xe tuần tra ít nhất trên đường khi đến trạm } i
  \text{ tại thời điểm } T.
\]
Công thức chuyển và điều kiện biên là
\[
  F(i,T)=\min_{300\le T_k\le600}
  \{F(i-1,T-T_k)+w(i-1,T-T_k,T)\}\qquad (2\le i\le n),
\]
\[
  F(1,\text{06:00:00})=0.
\]
Nghiệm bài toán là
\[
  \min\{F(n,T)\}.
\]
Trong đó, \(w(i-1,T-T_k,T)\) là số xe tuần tra mà xe mục tiêu gặp khi xuất
phát từ trạm \(i-1\) tại thời điểm \(T-T_k\) và đến trạm \(i\) tại thời điểm
\(T\).

Phân tích độ phức tạp thời gian: số giai đoạn của bài là \(n\), số trạng thái
ở giai đoạn \(i\) là \((i-1)\cdot300\). Vì vậy tổng số trạng thái là
\[
  O\left(\sum_{i=1}^{n}(i-1)\cdot300\right)
  =O\left(300\cdot\frac{n(n-1)}{2}\right)
  =O(150n^2).
\]
Mỗi trạng thái chuyển từ 300 trạng thái. Thời gian mỗi lần chuyển phụ thuộc
chủ yếu vào cách tính hàm \(w\). Dưới đây so sánh độ phức tạp khi dùng các
cách tính khác nhau.

\paragraph{Tính hàm \(w\).}
\term{Phương pháp 1.} Mỗi quyết định đều tính một lần, kiểm tra mọi xe tuần
tra xuất phát từ trạm \(i\). Như vậy thời gian trung bình mỗi lần chuyển là
\(O(1+m/n)\). Do \(m\) lớn nhất là 300, tổng thời gian là
\[
  O\left(150n^2\cdot300\cdot\left(1+\frac{m}{n}\right)\right)
  =O\left(\frac{m^2n^2}{2}+\frac{m^3n}{2}\right)
  =O(m^3n).
\]

\term{Phương pháp 2.} Quan sát kỹ công thức chuyển trạng thái, khi chuyển đến
trạng thái \(F(i,T)\), các hàm \(w\) được tính đều xuất phát từ trạm \(i\) và
cùng thời điểm xuất phát \(T\); chỉ thời điểm đến tương ứng là khác nhau. Ta
xét liệu có thể tìm quan hệ giữa chúng để tận dụng kết quả đã có và giảm tính
lặp hay không.

Xét quan hệ giữa \(w(i,T,k)\) và \(w(i,T,k+1)\). Với mỗi xe tuần tra xuất
phát từ trạm \(i\), giả sử thời điểm xuất phát và thời điểm đến của nó lần
lượt là \(\operatorname{Stime}\) và \(\operatorname{Ttime}\). Khi đó:
\begin{itemize}
  \item Nếu \(\operatorname{Ttime}<k\) hoặc \(\operatorname{Ttime}>k+1\), xe
  mục tiêu A và xe mục tiêu B có tình huống gặp xe tuần tra giống nhau.
  \item Nếu \(\operatorname{Ttime}=k\), xe mục tiêu A gặp xe tuần tra này.
  Với xe mục tiêu B: nếu \(\operatorname{Stime}\le T\), B không gặp xe tuần
  tra; ngược lại B cũng gặp.
  \item Nếu \(\operatorname{Ttime}=k+1\), xe mục tiêu B gặp xe tuần tra này.
  Với xe mục tiêu A: nếu \(\operatorname{Stime}\ge T\), A không gặp xe tuần
  tra; ngược lại A cũng gặp.
\end{itemize}

Đặt
\[
  \Delta[k]=w[i,T,k+1]-w[i,T,k].
\]
Từ thảo luận trên,
\[
  \Delta[k]=G((\operatorname{Ttime}=k+1)\wedge(\operatorname{Stime}\ge T))
  -G((\operatorname{Ttime}=k)\wedge(\operatorname{Stime}\le T)).
\]
Trong đó \(G(P)\) là số xe tuần tra xuất phát từ trạm \(i\) và thỏa điều kiện
\(P\).

Như vậy ta đã tìm được quan hệ giữa các giá trị của hàm \(w\). Khi chuyển
đến trạng thái \(F(i,T)\), trước hết quét một lần mọi xe tuần tra xuất phát
từ trạm \(i\); trong khi tính \(w[i,T,T+300]\), đồng thời tính
\(\Delta[T+301..T+600]\). Bước này mất \(O(m/n)\). Trong các chuyển trạng thái
sau đó, từ
\[
  w[i,T,k+1]=w[i,T,k]+\Delta[k],
\]
chỉ cần \(O(1)\) là tính được giá trị hàm \(w\), nên thời gian chuyển trạng
thái chỉ còn \(O(1)\). Tổng thời gian của thuật toán là
\[
  O\left(150n^2\left(\frac{m}{n}+300\right)\right)
  =O\left(\frac{m^2n^2}{2}+\frac{m^2n}{2}\right)
  =O(m^2n^2).
\]

Tuy bậc độ phức tạp thời gian của thuật toán không giảm, do \(m\le300\) và
\(n\le50\), hiệu quả tối ưu trong kiểm thử thực tế vẫn rất rõ rệt.

Tối ưu hóa bài này thực chất là áp dụng chính tư tưởng quy hoạch động. Khi
tính hạng hằng của công thức truy hồi, ta đưa vào hàm \(\Delta\), tận dụng
kết quả tính trước đó, tránh tính lặp, loại bỏ phần dư thừa và nâng cao hiệu
suất thời gian. Việc tính hàm \(w\) trong bài Post Office ở trên cũng dùng
tiền xử lý để giảm tính lặp. Gần đây còn xuất hiện quy hoạch động kép, cũng
áp dụng tư tưởng này để dùng quy hoạch động tính hạng hằng của công thức
truy hồi. Có thể thấy phương pháp tối ưu này có tính phổ quát cao.

\section{Kết luận}

Bài viết chủ yếu thảo luận tối ưu hóa hiệu suất thời gian của quy hoạch động
từ ba mặt: giảm tổng số trạng thái, giảm số trạng thái chuyển của mỗi trạng
thái, và giảm thời gian chuyển trạng thái. Đồng thời, bài viết cũng gián tiếp
thảo luận phương pháp tối ưu hóa thuật toán nói chung.

Trong quá trình tối ưu, tôi nhận ra rằng tối ưu thuật toán một mặt cần phân
tích sâu thuộc tính của bài toán và đào sâu bản chất của nó; mặt khác cần bắt
đầu từ chỗ còn yếu của thuật toán hiện có, liên tục cải tiến và theo đuổi sự
hoàn thiện.

Thiết kế thuật toán quy hoạch động có tính linh hoạt rất lớn, cần phân tích
từng mô hình cụ thể theo tình huống cụ thể. Thiết kế thuật toán là như vậy,
và tối ưu thuật toán cũng như vậy. Những điều trình bày trong bài chỉ là một
số phương pháp mang tính chung; nhiều kỹ thuật tối ưu còn cần các thí sinh đào
sâu trong huấn luyện và thi đấu hằng ngày.

Quy hoạch động, với tư cách là một loại thuật toán hiệu quả, vẫn còn nhiều
dư địa để tối ưu. Không ngừng nâng cao hiệu năng thuật toán để nó thích ứng
với quy mô lớn hơn, tôi nghĩ đó là mong muốn chung của đông đảo thí sinh tin
học. Hy vọng mọi người cùng nghiên cứu và thảo luận việc tối ưu hóa thuật
toán quy hoạch động; đó cũng là ý định ban đầu khi viết bài này.

\section*{Tài Liệu Tham Khảo}

\begin{enumerate}
  \item Wu Wenhu, Zhao Peng. \textit{Problems and analyses from the 1993--1996
  American computer programming contests}. Tsinghua University Press, 1999.
  \item Wu Wenhu, Wang Jiande. \textit{Analyses of international and domestic
  youth informatics contest problems}. Tsinghua University Press, 1997.
  \item Fu Qingxiang, Wang Xiaodong. \textit{Algorithms and data structures}.
  Publishing House of Electronics Industry, 1998.
  \item Organizing Committee of the National Youth Informatics Olympiad
  Regional League. \textit{Informatics Olympiad Quarterly}, 1999.3, 2000.2.
\end{enumerate}

\section*{Phụ Lục}

\begin{enumerate}
  \item Công thức mô tả độ phức tạp thời gian của quy hoạch động ở trên chỉ
  là mô tả trực quan các yếu tố quyết định độ phức tạp thời gian, không thể
  dùng như một công thức tính phổ quát.
  \item Bất đẳng thức tứ giác do Donald E. Knuth nêu ra từ cấu trúc dữ liệu
  cây tìm kiếm nhị phân tối ưu; ở đây nó được dùng để chứng minh tính đơn điệu
  của quyết định trong quy hoạch động.
  \item Dùng cây tra cứu để tìm chuỗi chỉ cần đi từ gốc đến nút lá, thời gian
  cần thiết tỉ lệ với độ dài chuỗi. Nếu hàm băm thật sự ngẫu nhiên, giá trị
  của hàm băm liên quan đến từng chữ cái trong chuỗi. Vì vậy, thời gian tính
  giá trị hàm băm xấp xỉ thời gian thực hiện một thao tác của cây tra cứu.
  Nhưng sau khi tính giá trị hàm băm còn phải xử lý va chạm. Do đó trong
  trường hợp thông thường, khi tra cứu chuỗi, cây tra cứu tiết kiệm thời gian
  hơn bảng băm.
\end{enumerate}

\section*{Phần Trình Chiếu Trích Xuất}

Phần sau đây là nội dung trình chiếu được trích xuất từ cuối PDF.

\subsection*{Tối ưu hóa hiệu suất thời gian của thuật toán quy hoạch động}

Mao Ziqing, Trường trung học số 3 Phúc Châu.

\[
\begin{aligned}
\text{độ phức tạp thời gian của quy hoạch động}
&= \text{tổng số trạng thái}\\
&\quad \times \text{số trạng thái chuyển của mỗi trạng thái}\\
&\quad \times \text{thời gian cho mỗi lần chuyển trạng thái}.
\end{aligned}
\]

\begin{enumerate}
  \item Giảm tổng số trạng thái:
  \begin{itemize}
    \item cải tiến biểu diễn trạng thái;
    \item các phương pháp khác, chẳng hạn chọn hướng quy hoạch thích hợp.
  \end{itemize}
  \item Giảm số trạng thái chuyển của mỗi trạng thái:
  \begin{itemize}
    \item giảm lượng quyết định theo tính chất của nghiệm tối ưu;
    \item các phương pháp khác, chẳng hạn dùng bất đẳng thức tứ giác để chứng
    minh tính đơn điệu của quyết định.
  \end{itemize}
  \item Giảm thời gian chuyển trạng thái:
  \begin{itemize}
    \item giảm thời gian quyết định bằng cấu trúc dữ liệu thích hợp;
    \item giảm thời gian tính công thức truy hồi bằng tiền xử lý, tận dụng kết
    quả đã tính, v.v.
  \end{itemize}
\end{enumerate}

\subsection*{Ví dụ trình chiếu 1: Raucous Rockers}

Có \(n\) bài hát của nhóm Raucous Rockers. Cần chọn một số bài để phát hành
trên \(m\) đĩa, mỗi đĩa tối đa \(t\) phút nhạc. Các bài không được chồng lấn,
bộ đĩa phải theo thứ tự sáng tác, và tổng số bài chứa được là nhiều nhất. Cho
\(n,m,t\) và độ dài \(n\) bài hát theo thứ tự sáng tác; không bài nào dài hơn
một đĩa, và không thể đặt tất cả bài lên đĩa. Hãy xuất số bài nhiều nhất.

Với độ dài các bài là \(\operatorname{long}[1..n]\), trạng thái ban đầu:
\[
  g[i,j,k]\quad (0\le i\le n,\ 0\le j\le m,\ 0\le k<t),
\]
biểu thị số bài hát nhiều nhất có thể ghi khi xét \(i\) bài đầu, dùng \(j\)
đĩa cộng thêm \(k\) phút. Khi \(k\ge\operatorname{long}[i]\) và \(i\ge1\):
\[
  g[i,j,k]=\max\{g[i-1,j,k-\operatorname{long}[i]]+1,\ g[i-1,j,k]\}.
\]
Khi \(k<\operatorname{long}[i]\) và \(i\ge1\):
\[
  g[i,j,k]=
  \max\{g[i-1,j-1,t-\operatorname{long}[i]]+1,\ g[i-1,j,k]\}.
\]
Điều kiện biên:
\[
  g[0,j,k]=0\qquad (0\le j\le m,\ 0\le k<t).
\]
Nghiệm tối ưu là \(g[n,m,0]\), thời gian \(O(nmt)\).

Trạng thái cải tiến:
\[
  g[i,j]=(a,b),
\]
biểu thị khi chọn \(j\) bài trong \(i\) bài đầu, số đĩa ít nhất cần dùng là
\(a\) đĩa cộng thêm \(b\) phút. Công thức:
\[
  g[i,j]=\min\{g[i-1,j],\ g[i-1,j-1]+\operatorname{long}[i]\}.
\]
Nếu \((a,b)+\operatorname{long}[i]=(a',b')\), thì
\[
\begin{array}{ll}
b+\operatorname{long}[i]\le t: & a'=a,\quad b'=b+\operatorname{long}[i],\\
b+\operatorname{long}[i]> t: & a'=a+1,\quad b'=\operatorname{long}[i].
\end{array}
\]
Điều kiện biên là \(g[i,0]=(0,0)\), \(0\le i\le n\). Đáp án:
\[
  \operatorname{answer}=\max\{k\mid g[n,k]\le (m-1,t)\}.
\]
Thời gian là \(O(n^2)\).

\subsection*{Ví dụ trình chiếu 2: điểm lớn nhất trong bài gộp đá}

Trên một vòng tròn có \(n\) đống đá. Cần gộp có thứ tự thành một đống; mỗi lần
chỉ được gộp hai đống kề nhau, và số đá của đống mới là điểm của lần gộp. Hãy
tìm điểm nhỏ nhất, điểm lớn nhất và phương án tương ứng. Ở đây chỉ xét điểm
lớn nhất.

Gọi số đá là \(d[1..n]\). Đặt
\[
  t[i,j]=\sum_{k=i}^{j}d[k].
\]
Trạng thái:
\[
  m[i,j]\quad (1\le i,j\le n),
\]
biểu thị điểm lớn nhất khi gộp \(d[i..j]\). Công thức:
\[
  m[i,j]=\max_{i\le k\le j-1}\{m[i,k]+m[k+1,j]+t[i,j]\}\qquad (i<j),
\]
với điều kiện biên \(m[i,i]=0\). Đặt \(s[i,j]=k\), biểu thị vị trí tách tối
ưu. Thời gian là \(O(n^3)\).

Vị trí tách tối ưu khi gộp đống \(i\) đến đống \(j\), tức \(s[i,j]\), hoặc
bằng \(i\), hoặc bằng \(j-1\). Nói cách khác, phương án tối ưu chỉ có thể là
\[
  \{(i)(i+1\ldots j)\}\quad\text{hoặc}\quad
  \{(i\ldots j-1)(j)\}.
\]
Chứng minh phác thảo trong phần trình chiếu dùng phản chứng. Giả sử
\(s[i,j]=p\) và \(i<p<j-1\). Nếu \(t[i,p]\le t[p+1,j]\), đặt \(q=s[i,p]\).
Phương án tối ưu đang xét là
\[
  \{[(i\ldots q)(q+1\ldots p)](p+1\ldots j)\},
\]
có điểm
\[
  F_1=m[i,q]+m[q+1,p]+m[p+1,j]+t[i,j]+t[i,p].
\]
Nhưng có thể dựng phương án
\[
  \{(i\ldots q)[(q+1\ldots p)(p+1\ldots j)]\},
\]
có điểm
\[
  F_2=m[i,q]+m[q+1,p]+m[p+1,j]+t[i,j]+t[q+1,j].
\]
Vì \(q<p\), ta có \(t[i,p]\le t[p+1,j]<t[q+1,j]\), nên \(F_1<F_2\), mâu
thuẫn với giả thiết tối ưu. Trường hợp còn lại tương tự.

Công thức chuyển tối ưu:
\[
  m[i,j]=\max\{m[i+1,j],m[i,j-1]\}+t[i,j]\qquad (i<j),
\]
với \(m[i,i]=0\). Thời gian là \(O(n^2)\).

\subsection*{Ví dụ trình chiếu 3: LOSTCITY}

Cho một bảng từ, quy tắc ngữ pháp xác định và một bài văn. Bài văn và bảng từ
chỉ chứa 26 chữ cái thường tiếng Anh \(a\ldots z\). Từ trong bảng từ chỉ có
danh từ, động từ và phó từ; các từ cùng từ loại đôi một khác nhau. Số từ không
vượt quá 1000, độ dài mỗi từ không vượt quá 20. Quy tắc ngữ pháp: cụm danh từ
là một số phó từ tiền tố nối với một danh từ; cụm động từ là một số phó từ
tiền tố nối với một động từ; câu bắt đầu bằng cụm danh từ, sau đó cụm danh từ
và cụm động từ nối xen kẽ. Bài văn dài không quá 5000 ký tự. Biết rằng bài văn
do hữu hạn câu tạo thành, mỗi câu có hữu hạn từ. Cần chia bài văn thành ít
câu nhất; trong điều kiện đó, số từ được chia cũng ít nhất.

Lần lượt dùng \(v,u,a\) để biểu thị động từ, danh từ và phó từ; bài văn là
\(L[1..M]\). Trạng thái:
\[
\begin{array}{ll}
F(v,i): & \text{phương án tối ưu khi } i \text{ ký tự đầu kết thúc bằng động từ,}\\
F(u,i): & \text{phương án tối ưu khi } i \text{ ký tự đầu kết thúc bằng danh từ.}
\end{array}
\]
Khi \(i\ne M\), cuối trạng thái có thể kèm một số phó từ hậu tố. Công thức:
\[
\begin{aligned}
F(v,i)=\min\{&
F(u,j)+(0,1), && L(j+1..i)\text{ là động từ};\\
&F(v,j)+(0,1), && L(j+1..i)\text{ là phó từ},\ i\ne M\},
\end{aligned}
\]
\[
\begin{aligned}
F(u,i)=\min\{&
F(u,j)+(1,1), && L(j+1..i)\text{ là danh từ};\\
&F(v,j)+(0,1), && L(j+1..i)\text{ là danh từ};\\
&F(u,j)+(0,1), && L(j+1..i)\text{ là phó từ},\ i\ne M\}.
\end{aligned}
\]
Điều kiện biên:
\[
  F(v,0)=(1,0),\qquad F(u,0)=(\infty,\infty).
\]
Đáp án:
\[
  \min\{F(v,M),F(u,M)\}.
\]

So sánh các phương pháp tra chuỗi, với \(N\le1000\) là số từ và \(M\le5000\)
là độ dài bài văn:
\[
\begin{array}{lcccc}
\hline
& \text{tìm tuần tự} & \text{tìm nhị phân} & \text{bảng băm} & \text{cây tra cứu}\\
\hline
\text{thời gian} & O(20MN) & O(20M\log_2N) & O(20M) & O(M)\\
\text{số so sánh xấu nhất} & 10^8 & 10^6 & 10^5 & 5\cdot10^3\\
\text{Input.009} & \text{quá giờ} & 1.27\text{s} & 0.32\text{s} & 0.05\text{s}\\
\text{Input.010} & \text{quá giờ} & 1.33\text{s} & 0.33\text{s} & 0.05\text{s}\\
\hline
\end{array}
\]
Môi trường kiểm thử trong phần trình chiếu: Pentium 200, 32 MB bộ nhớ.

\end{document}
